\section{Step 15 Form each candidate from the plane origin}\label{h263-motion-s05}
For every remaining admitted $(d_x,d_y)$, form $(x+d_x,y+d_y)$ first, then its
byte offset using the reference stride. Add that offset to the reference plane
pointer. Recomputing from the origin avoids carrying a pointer from the end of
one candidate's sixteen rows into the next candidate.

For the worked position $(16,32)$ and displacement $(-2,-2)$, the candidate begins
at $(14,30)$. With stride 176, its offset is $30\times176+14=5294$ bytes from the
reference origin. Its last row begins another $15\times176=2640$ bytes later.
The current pointer continues to identify the original current block.

\paragraph{Observe} A load from the candidate's first row reads reference
samples 14 through 29 of row 30. It reads no column from the following row.
\Needspace{8\baselineskip}
\paragraph{Check} Which pointer changes when selecting a new candidate?
\begin{itemize}
\item[A.] The reference window pointer derived from the plane origin.
\item[B.] Both image pointers by the same displacement.
\item[C.] The result pointer by sixteen rows.
\end{itemize}
